% ============================================================
%  第五部分：工程映射（三级体系机制剖析）
% ============================================================

\section{工程映射：三级 Regime 体系的架构与机制剖析}
\label{sec:engineering}

前面七章搭建了地图，本章把它落到一类具体的工程形态上：一个三级架构的期货 Regime 识别体系——\textbf{L1 全市场判档}（七维量化 $+$ 五状态）、\textbf{L2 产业链层}（HMM 三窗口投票 $+$ 规则法）、\textbf{L3 品种层}（三信号规则法 $+$ 产业链条件化），配套\textbf{多源数据降级链}与\textbf{ML 因子模块}。以下按“先架构、后机制”的顺序展开，聚焦可迁移的设计逻辑与参数结构，不指向任何具体实现。

\subsection{体系总览：分辨率分级的三个问题}

\begin{table}[htbp]
\centering
\footnotesize
\begin{tabular}{>{\raggedright\arraybackslash}p{2.1cm}>{\raggedright\arraybackslash}p{3.0cm}>{\raggedright\arraybackslash}p{3.7cm}>{\raggedright\arraybackslash}p{2.9cm}}
\toprule
\textbf{层级} & \textbf{回答的问题} & \textbf{方法} & \textbf{输出} \\
\midrule
L1 全市场 & 用多大仓位、什么姿态？ & 七维量化判档．五状态 & 下游约束参数 \\
 & & & （杠杆/权重/止损） \\
L2 产业链 & 钱放在哪条链？ & HMM 三窗口投票 $+$ 规则法 & 状态 $+$ 概率 $+$ 驱动力 \\
 & & & $+$ 分散度 \\
L3 品种 & 具体做哪个品种？ & 三信号规则法 $+$ 条件化 & 品种状态 $+$ 置信度 \\
\bottomrule
\end{tabular}
\caption{三级架构的分工：每一级都不假装回答不属于它的问题}
\label{tab:three-levels}
\end{table}

这个设计对应了第一章“状态是决策分辨率的函数”的观点：全市场的状态决定系统性暴露，产业链的状态决定截面配置，品种的状态决定执行层面的取舍。三级之间通过两个接口连接：L1 向下输出\emph{姿态约束}（参数边界，而非具体仓位）；L2 的\emph{状态概率}作为 L3 的\emph{条件化先验}。

\textbf{三级接口流（数据如何逐级传导）}：
\begin{center}
\footnotesize
\renewcommand{\arraystretch}{1.4}
\begin{tabular}{@{}p{2.2cm}>{\raggedright\arraybackslash}p{4.4cm}>{\raggedright\arraybackslash}p{4.4cm}@{}}
\toprule
级 & 输入 & 输出向下传导 \\
\midrule
L1 & 七维全市场因子 & 姿态约束（杠杆上限/权重边界/止损开关）$\to$ 下游优化器 \\
L2 & 9 条产业链等权指数 & 状态概率 $+$ 驱动力 $\to$ 作为 L3 条件化先验 \\
L3 & 品种三信号 $+$ L2 先验 & 品种状态 $+$ 置信度 $\to$ 具体合约执行 \\
\bottomrule
\end{tabular}
\end{center}
\textbf{两级接口的容错设计}：L1 只给"参数边界"不给"具体仓位"（第 \ref{sec:engineering} 节 L1 约束协议）；L2 概率经 $1.15/0.85$ 放大/收缩后进 L3，而非直接指令。这样上游识别误差被约束吸收，不会沿链放大成错误仓位——\textbf{三层不是简单的"级联漏斗"，而是一个带接口缓冲的层级结构}。

\subsection{L1：全市场七维判档体系}

\subsubsection{为什么是七个维度}

L1 的维度设计从期货市场的结构特性推导，而非照搬股票市场的“广度/情绪/盈利”框架。股票状态的核心是单一市场（Beta）与风格（成长/价值）；期货是多资产截面（Alpha）$+$ 期限结构 $+$ 基本面库存三重结构。因此七维为：

\begin{enumerate}[itemsep=2pt]
\item \textbf{跨资产趋势结构}：各板块代表品种的 ADX(14) 中位数、板块动量一致性、趋势品种占比——回答“是否存在系统性趋势”；
\item \textbf{期限结构状态}：全市场展期收益率中位数、Backwardation/Contango 占比、斜率变化——回答“Carry 是否可用、趋势是否有基本面支撑”；
\item \textbf{波动率环境}：全市场波动率中位数及其历史分位——回答“杠杆空间与风险预算”；
\item \textbf{资金与情绪}：持仓量变化、龙虎榜多空集中度、期权隐含波动率——回答“拥挤度与参与者预期”；
\item \textbf{宏观与美元环境}：美元指数、美债收益率、中美利差、PMI、通胀——回答“宏观驱动是否存在”；
\item \textbf{库存与供需周期}：去库/累库品种占比、库存与价格方向的一致性——回答“基本面与价格是否共振”；
\item \textbf{ML 因子有效性}：滚动 IC 均值/标准差、特征重要性稳定性、早停轮数——回答“市场是否存在可预测的非线性结构”。
\end{enumerate}

\begin{remark}[第七条维度的设计哲学]
前六维都是“市场自身的量价与基本面描述”，第七维是一个\emph{元信号}：它不描述市场，而描述“我们的模型在当前市场上是否有效”。其背后的逻辑是：若 ML 因子滚动 IC 持续为正且特征重要性稳定，说明价格序列中存在可预测结构（趋势友好的证据）；若 IC 在零附近波动，说明市场接近随机（均值回归友好的证据）；若 IC 标准差急剧放大，说明结构不稳定（切换/混沌的证据）。这种“以模型表现反推市场状态”的元信号设计，是对常规量价维度的正交补充，也是本体系中少见的原创性设计。
\end{remark}

\subsubsection{五状态判档与下游约束协议}

七维证据聚合为五个状态：\emph{趋势友好型、均值回归友好型、宏观驱动型、高波动防御型、切换/混沌型}。每个状态附带明确的维持条件、切换前置信号与失效标志（例如：趋势友好的失效标志为“指数跌破关键均线、持仓量萎缩”；高波动防御的维持条件是“波动率 $> 80\%$ 分位”）。

关键工程接口是把状态\emph{翻译成数字约束}，直接限定下游系统：

\begin{table}[htbp]
\centering
\footnotesize
\begin{tabular}{p{2.6cm}p{1.4cm}p{1.9cm}p{1.4cm}p{1.7cm}p{1.7cm}}
\toprule
\textbf{约束参数} & \textbf{趋势友好} & \textbf{均值回归友好} & \textbf{宏观驱动} & \textbf{高波动防御} & \textbf{切换/混沌} \\
\midrule
板块权重上限 & 30\% & 20\% & 25\% & 15\% & 15\% \\
总杠杆上限 & 2.0x & 1.5x & 1.8x & 1.0x & 1.0x \\
单品种仓位上限 & 8\% & 6\% & 7\% & 4\% & 4\% \\
止损收紧倍率 & 1.0x & 1.0x & 1.0x & 0.7x & 0.8x \\
ML 因子权重上限 & 30\% & 5\% & 20\% & 10\% & 10\% \\
\bottomrule
\end{tabular}
\caption{L1 状态 $\to$ 下游约束参数（节选自体系的产出协议）}
\label{tab:l1-constraints}
\end{table}

\begin{remark}[为什么输出“约束”而不是“仓位”]
这是一个值得强调的设计决策：L1 不直接给出仓位数字，而是给出\emph{参数边界}（上限、倍率、开关）。理由是分工与稳健性——L1 的状态判定误差不会直接变成错误仓位，而是在下游优化器内部被约束吸收；同时下游（因子系统）保留自身逻辑的完整性，只被“调参”而不会被“夺权”。这种“约束式接口”比“指令式接口”对上游识别误差的容错性高一个量级。
\end{remark}

\subsection{L2：产业链层——HMM 三窗口投票机制}

L2 对每条产业链（体系中配置了 9 条：黑色、有色、贵金属、新能源、能源化工、农产品、金融期货、纸浆造纸、橡胶）合成等权指数并识别状态。以下按执行顺序剖析各环节的机制设计。

\subsubsection{合成指数}

合成指数环节以首个品种的日期轴为基准，对齐成员收盘价（前向 $+$ 后向填充），做“首日归一化 $+$ 等权平均”得到产业链指数，成交量取成员之和。设计取舍清晰：等权而非市值加权，避免了单一权重品种绑架产业链状态；归一化使跨品种的量纲差异消解。

\subsubsection{HMM 三窗口投票机制}

状态识别的核心是“三窗口投票”：对 $63$ / $126$ / $252$ 天三个窗口分别训练一次 HMM，各投一票，取得票最多的状态。单窗口的处理流程如下（伪代码）：

\begin{lstlisting}[language={}, caption={单窗口 HMM 识别流程（伪代码）}]
输入: 产业链指数收益率序列, 窗口长度
输出: 窗口投票的状态标签

特征 <- [日收益率, 20日滚动波动率, 20日振幅]
特征 <- 标准化(特征)
最优模型 <- 空
重复 5 次:
    候选 <- 训练高斯HMM(状态数=5, 全协方差, 迭代上限 50)
    若 候选对数似然 > 最优模型对数似然: 最优模型 <- 候选
状态序列 <- 维特比解码(最优模型, 特征)
当前隐状态 <- 状态序列末位
按该状态画像(平均收益 / 平均波动)语义映射 -> 牛/熊/震荡/高波/低波
\end{lstlisting}

这部分有三个值得评述的设计点：

\begin{enumerate}[itemsep=3pt]
\item \textbf{三窗口投票 $=$ 多尺度证据的硬表决}。$63$/$126$/$252$ 天覆盖季度、半年、年度三个尺度。它的工程目标是\emph{抗单尺度偶然性}：任何单一窗口拟合出的“伪状态”都难以在三个尺度上同时成立。这与第三部分“投票是工作点编码”的框架完全一致——它换取的正是误报的下降，代价是切换确认的延迟。
\item \textbf{状态映射是“画像解释”而非“原样输出”}。HMM 训练出的 5 个隐状态本身没有语义（第二章的 label switching 问题）；识别流程用\emph{当前状态画像的}平均收益与平均波动（标准化尺度）做二次映射：平均波动或平均收益超过相应阈值（$0.8$ / $0.5$）才标记为高波/牛/熊，否则归为震荡。这正是“以状态画像为身份”的做法在工程中的朴素版本。
\item \textbf{降级路径的完备}。若统计库缺失或某窗口样本不足（$< 30$ 个特征点），该窗口自动降级为单窗口规则法投票，保证“永远有输出”。这在生产上是正确取舍——但第三部分的检查清单提醒：降级分支的输出置信度语义与 HMM 分支不同，系统应在报告中显式标注（当前设计由“置信度分”字段部分承担这一职责）。
\end{enumerate}

\subsubsection{规则法与合成逻辑}

规则法分支用收益率/波动率的 z-score 与均线排列（20 日 $/$ 60 日均线）产生状态，再与 HMM 结果合成：\emph{两者一致时置信度加成}（均值 $+ 0.1$，上限 $0.95$）；\emph{不一致时按 HMM $0.6$ 对规则 $0.4$ 的加权裁决}。这是一个典型的分工：统计模型提供主体判断（权重更高），规则法提供约束与可解释兜底。

\begin{remark}[两个必须如实记录的实现细节]
第一，对外输出的“状态概率”字段由启发式加权生成（基础概率 $0.1$ 加 HMM 命中项与规则命中项后归一化），\textbf{并非 HMM 的真实后验概率}——真实后验概率接口并未参与计算。下游若把该字段当作贝叶斯后验使用（例如做期望效用加权），会引入系统性偏差。第二，反馈到“当前状态”判定的解码方式是 Viterbi 路径（路径级联合最优），与第~\ref{sec:statespace}~节讨论的“逐点边缘最优”存在已知差异。这两点不影响系统作为“判别器”可用，但构成明确的接口语义债务，建议在文档或字段命名中显式声明。
\end{remark}

\subsubsection{驱动力识别与分散度}

驱动力识别环节依据“上游 $-$ 下游趋势差”与原油品种存在性，将状态归因为原油驱动/成本驱动/需求驱动/宏观驱动；分散度环节以成员品种 20 日趋势的\emph{变异系数}度量产业链内部分散度（$0$ 为一致、$1$ 为高度分化）。两个指标都以“可解释、可审计”为优先，放弃了对假设的严格性——与第一章的“判别式刻画”定位一致。

\subsection{L3：品种层——三信号规则法}

\subsubsection{因子计算层}

因子计算层产出六组因子：多周期动量（5/20/60 日 $+$ 加速度）、波动率（20 日年化 $+$ 历史分位 $+$ 波幅占比）、均线（快慢交叉信号、价格偏离）、ADX/DI（趋势强度与方向）、量价（量比、量价同步性）、位置（60 日价格位置、20 日突破度）。随后合成三个信号：

\begin{align}
\text{趋势投票}\;&=\;\mathrm{clip}\Big(\frac{\sum_i w_i\, v_i}{\sum_i w_i},\; -1,\; 1\Big), \quad
\text{其中动量 } v = \mathrm{clip}(r \times 10),\ w = 1/\text{窗口},\\
&\qquad\text{均线交叉 } w = 2.0,\quad \text{DI 差 } w = 1.5,\quad \text{量价 } w = 0.5;\\
\text{波动信号}\;&=\;\text{20 日波动率的历史分位 } \in [0,1];\\
\text{OI 信号}\;&=\;\text{成交量变化率的 z 值}.
\end{align}

\begin{remark}[OI 代理：一个必须标注的工程妥协]
设计文档明确记录：OI 信号因数据源（免费版接口）不支持持仓量数据，故以成交量变化作为代理。这是一个诚实且合理的降级——但语义上会引入偏差：成交量放大既可能来自持仓增长（新资金入场），也可能来自存量换手（无方向性信息）。体系另有一条更优路径：本地数据库的主力连续表自带持仓量列，真 OI 可用——切换数据源即自动升级信号质量。这种“同一信号在降级链不同层级有不同精度”的问题，是所有多源系统的共有课题。
\end{remark}

\subsubsection{状态判定与置信度}

状态判定采用如下判定树：\emph{流动性检查}（量比 $< 0.3$ $\to$ 流动性枯竭）$\to$ \emph{高波动检查}（波动分位 $\ge 0.8$ 且非“强趋势 $+$ ADX 确认”，判高波动；若是强趋势则维持趋势判定——注意这个分支体现了“高波动中的强趋势是趋势而非噪声”的判断）$\to$ \emph{趋势分级}（$|$趋势投票$| \ge 0.5$ 强趋势 / $\ge 0.25$ 弱趋势 / 否则均值回归）。置信度由基础分（$0.5 + 0.3\times|$趋势投票$|$）经 ADX 确认（$+0.1$）、OI 同向（$+0.05$）或背离（$-0.05$）修正后裁剪到 $[0.3, 0.95]$。

\subsubsection{产业链条件化先验}

产业链条件化环节把 L2 状态作为 L3 的先验：顺链（产业链牛 $+$ 品种涨）放大置信度 $1.15$ 倍，逆链收缩至 $0.85$；高波共振 $1.10$、震荡共振 $1.05$；若 L2 状态概率的最大值 $>0.7$（产业链判断高确定性），再整体乘 $1.05$。

\begin{remark}[条件化是“启发式贝叶斯”，不是贝叶斯]
严格的条件化应计算 $\Prob(\text{品种状态} \mid \text{产业链状态})$，即用产业链状态作为先验更新品种状态的后验。当前实现是这一目标的\textbf{近似}：用固定比例的乘性因子模拟“顺链增强、逆链削弱”。其优点（透明、无需联合标注、参数可审）与缺点（缩放幅度是先验设定，非数据估计）同样明显。最省力的改进方向是：用历史数据估计“L2 状态与 L3 状态的条件一致率”，据此校准缩放系数——保留启发式框架，把拍定的常数交给数据。
\end{remark}

\begin{remark}[三个“配置声明未消费”的工程发现]
逐项核对配置声明与判定逻辑，发现三处“声明了但未被消费”的参数，如实记录：（1）OI 确认阈值（$1.0$）声明后未在判定逻辑中使用（OI 确认分支的实际条件等价于“OI 信号为正”）；（2）最小成交量门槛（$20000$）未使用（流动性检查实际采用量比 $< 0.3$）；（3）HMM 参数段声明的“迭代上限 $100$、初始化 $10$ 次”未生效（实际执行为 $50$ 次迭代、$5$ 次初始化），收敛容差未传递。这不是缺陷指控，而是所有演进型工程的常态——但应当被治理：配置治理的最小方案是让“读配置的代码”与“配置文件”做一次性对账，把未消费项删除或接入，避免未来读者对名义参数产生错误预期。
\end{remark}

\subsection{数据工程：多源降级链与就绪度}

\subsubsection{降级链设计}

\begin{table}[htbp]
\centering
\footnotesize
\begin{tabular}{p{2.4cm}p{4.7cm}p{4.3cm}}
\toprule
\textbf{优先级} & \textbf{数据源与接口} & \textbf{特点与注意} \\
\midrule
1 & 本地客户端（本机接口，主力合约日线，$\sim$400 根） & 最新鲜；无 OI；主力合约无拼接历史 \\
2 & 第三方连续合约接口（需凭据） & 连续合约；免费版限速，逐品种连接 \\
3 & 本地数据库（主力连续表） & 全深度历史 $+$ \textbf{完整 OI}；依赖本地同步 \\
4 & 本地缓存（列式文件） & 离线兜底；字段可能不全（缺成交额自动补算） \\
\bottomrule
\end{tabular}
\caption{数据源自动降级链（自动模式按序尝试）}
\label{tab:data-chain}
\end{table}

\subsubsection{三个工程细节}

\begin{enumerate}[itemsep=3pt]
\item \textbf{合约解析}：以正则取合约代码的字母前缀（如 IC2609.CFF 归为 IC），忽略交易所后缀。该解析附带自检断言（IC2609.CFF $\to$ IC、AP2610.CZC $\to$ AP），把曾经的解析修复固化为回归测试——这是“缺陷驱动固化”的良好实践。
\item \textbf{成交额修正}：成交额按“成交量 $\times$ 三价均值 $\times$ 合约乘数”修正；乘数优先从合约信息表读取，缺失时用覆盖约 60 个品种的手工表兜底。乘数错误会以同样倍数污染所有成交额派生指标（如成交额分位、流动性筛选），这是跨市场因子工程中最容易静默出错的环节之一。
\item \textbf{就绪度门槛}：单品种 $\ge 120$ 行；产业链合成指数 $\ge 60$ 天；HMM 单窗口可用门槛为“窗口 $+$ 20”天（因为特征构造需 20 天滚动）。门槛以显式条件存在，且不满足时输出带原因的降级信息（例如“数据不足（少于 60 天）”）——这是就绪度检查的合格形态：不是静默返回空值，而是带着原因降级。
\end{enumerate}

\subsection{ML 因子模块：梯度提升树的滚动训练}

第七维的实现是一个自包含的 ML 管线（特征工程与模型训练两个模块）。其改进版相对初版的每一处修改都对应一个具体的过拟合教训：

\begin{itemize}[itemsep=2pt]
\item \textbf{时序验证集拆分}：从训练集末尾切出 20\% 作为验证集（\emph{按时间切，而非随机切}——避免未来信息泄漏），配合早停（20 轮）；
\item \textbf{容量与正则双降}：叶节点数 $31\to15$、树深 $6\to4$、叶子最小样本数 $50\to200$、L1/L2 正则系数 $0.1\to1.0$；
\item \textbf{集成}：3 个随机种子模型平均，降低单模型方差不稳定；
\item \textbf{诊断字段}：逐窗口记录训练/测试 IC（Spearman）、特征重要性 Top5、早停轮数——后者直接接入第七维：\emph{平均早停轮数过小、特征重要性排名剧烈波动，都是“信号弱 $+$ 结构不稳”的定量证据}。
\end{itemize}

滚动结构为“3 年训练 $/$ 1 年测试”。这套设计在方法论上完全对齐第~\ref{sec:validation}~节的清单（时序拆分、样本外评估、成本计入）；其输出的 IC 时间序列进入 L1 判档，使“模型有效性”成为判断市场状态的一项证据。

\subsection{跨资产推广：从商品期货到股债商汇}

体系天然包含跨资产种子：配置中的\emph{金融期货}产业链（股指期货 $+$ 国债期货）与\emph{贵金属}链。把架构推广到更宽的多资产谱系，原则是\textbf{保留三级分辨率与接口协议，替换每级的状态维度}——维度必须来自该资产类别的结构特性：

\begin{table}[htbp]
\centering
\footnotesize
\begin{tabular}{p{2.1cm}p{5.2cm}p{4.2cm}}
\toprule
\textbf{资产类} & \textbf{状态维度的“结构特性”替换} & \textbf{数据与注意事项} \\
\midrule
商品期货 & 期限结构、库存周期、开工率（现有体系） & 基本面维度依赖人工采集时效 \\
股指期货 & 基差/升贴水结构、成交额与宽度、两融与资金 & 市场广度需另接成分股数据 \\
国债期货 & 收益率曲线形态（陡平）、资金面（回购利率）、供给 & 曲线形态即“期限结构”同构物 \\
外汇 & 利差与政策差、离岸基差、央行干预信号 & 需引入汇率与离岸数据源 \\
\bottomrule
\end{tabular}
\caption{跨资产推广：同构架构、异质维度}
\label{tab:cross-asset}
\end{table}

\subsubsection{股票市场视角的对照}

同一架构在股票市场的既有实现（六维框架：指数趋势、市场广度、主线清晰度、情绪温度、宏观流动性、盈利预期；四状态：进攻/防守/震荡/切换）可以作为对照来检验“同构性”是否真实成立：

\begin{itemize}[itemsep=2pt]
\item \textbf{同构处}：三级分辨率（市场 $\to$ 板块/产业链 $\to$ 个股/品种）、状态画像化、切换信号可跟踪、输出约束而非指令——这些结构在股票版本中全部保留；
\item \textbf{异质处}：期货的状态维度围绕“期限结构/库存/展期收益”，股票围绕“广度/情绪/盈利预期”；期货的状态语义是“趋势/Carry/防御”，股票是“进攻/防守”——\emph{维度与语义都换了，架构与纪律没有换}。
\end{itemize}

\begin{remark}[推广的正确姿势]
跨资产推广最容易犯的错误是把来源市场的维度清单照搬到目标市场（例如把商品的“库存周期”维度硬塞给股指）。正确的姿势是回到第一部分的问题定义：为每个市场的状态识别重新回答“这个市场的非平稳性以什么方式表达”。架构是模板，维度是定制——这与“状态画像才是身份”的结论一脉相承。
\end{remark}

\keynote{三级架构的本质是把“市场状态”分解为三个分辨率的问题：全市场判档决定“用多大仓位、什么姿态”，产业链层决定“钱放在哪条链”，品种层决定“具体做哪个合约”。每一级都不假装回答不属于它的问题——这就是这套体系最大的工程美德，也是它可被跨资产复制的原因。}